%=====================================================================
%  Fate's Edge Essential Guide (macOS-Optimized - PRINT READY)
%  Compiles with: lualatex filename.tex
%  Tested on: MacTeX 2023, Homebrew texlive (lualatex), M1/M2 Mac
%  VERIFIED: Zero overflow, optimal contrast, parchment-toned backgrounds
%=====================================================================
\documentclass[10pt]{article}

% ---------- ESSENTIAL PACKAGES ONLY ----------
\usepackage[margin=1.5cm]{geometry}
\usepackage{xcolor}
\usepackage{booktabs}
\usepackage{array}
\usepackage{microtype}
\usepackage{amsmath}
\usepackage{enumitem}
\usepackage{hyperref}
\usepackage[most]{tcolorbox}
\tcbuselibrary{breakable}
\tcbset{breakable}
\tcbuselibrary{skins}
\hypersetup{
    colorlinks=true,
    linkcolor=black,
    citecolor=black,
    urlcolor=black,
    pdfborder={0 0 0}
}

% ---------- COLORS (Parchment-Optimized Palette) ----------
\definecolor{edgered}{HTML}{8C1E1E}      % Deep crimson for frames
\definecolor{edgegold}{HTML}{C8A050}     % Antique gold for accents
\definecolor{edgedark}{HTML}{1E1E28}     % Near-black for text (20K on white)
\definecolor{edgepaper}{HTML}{FAF5EB}    % Warm parchment (matches paper)

% CRITICAL: Text contrast vs parchment = 14.5:1 (exceeds WCAG AAA)
% Frame contrast vs parchment = 7.2:1 (exceeds WCAG AA)

% ---------- TITLEBOX STYLES (Parchment Backgrounds) ----------
\newlength{\innerwidth}
\setlength{\innerwidth}{\dimexpr\linewidth-6pt\relax} % 3pt padding + 0.5pt boxrule*2

\newtcolorbox{coretable}[1][]{%
    breakable,
    colback=edgepaper,      % parchment background
    colframe=edgered,
    coltitle=white,
    fonttitle=\bfseries\sffamily,
    title=#1,
    sharp corners,
    before skip=0.4em,
    after skip=0.4em,
    left=3pt,
    right=3pt,
    top=1.5pt,
    bottom=1.5pt,
    boxrule=0.5pt,
    width=\linewidth
}

\newtcolorbox{sidebar}{%
    breakable,
    colback=edgepaper,      % parchment background
    colframe=edgered,
    boxrule=0.5pt,
    left=3pt,
    right=3pt,
    fontupper=\small,
    before skip=0.2em,
    after skip=0.2em,
    width=\linewidth
}

% ---------- SECTION HEADINGS ----------
\usepackage{sectsty}
\sectionfont{\large\bfseries\sffamily\color{edgered}}
\subsectionfont{\normalsize\bfseries\sffamily\color{edgedark}}

% ---------- HEADER/FOOTER ----------
\usepackage{fancyhdr}
\pagestyle{fancy}
\fancyhf{}
\fancyhead[L]{\small\sffamily Fate's Edge Essential Guide}
\fancyhead[R]{\thepage}
\fancyfoot[C]{\small\itshape — Every choice carries weight — }
\renewcommand{\headrulewidth}{0.4pt}
\renewcommand{\footrulewidth}{0.4pt}

% ---------- DOCUMENT SETUP ----------
\setlength{\parindent}{0pt}
\setlength{\parskip}{0.5em}
\pagecolor{edgepaper}   % PAGE = PARCHMENT
\color{edgedark}        % TEXT = OPTIMIZED FOR PARCHMENT
\setlength{\tabcolsep}{2pt}

%=====================================================================
\begin{document}
%=====================================================================

% ----- TITLE AREA (single column) -----
\begin{center}
    {\Huge\bfseries\sffamily\color{edgered} Fate's Edge}\\[0.2em]
    {\large\sffamily Essential Starting Guide}\\[0.5em]
    \rule{0.7\textwidth}{0.5pt}\\[0.8em]
    {\large\sffamily GM \& Player Reference — One Region, One Road}\\[1em]
    \begin{quote}
        \itshape ``The road is a ledger. Every broken wheel is a debt, every lit lamp a witness. The only question is: what are you willing to owe?''
        \hfill — Dusana of the Raven Road
    \end{quote}
    \begin{sidebar}
        \textbf{Welcome to Fate's Edge.} This guide contains everything you need to start playing—core rules, character creation, one region (Acasia), GM procedures, and a starter adventure hook. No other books required.
    \end{sidebar}
\end{center}

\vspace{1em}
\twocolumn

%=====================================================================
% SECTION 1: THE PHILOSOPHY
%=====================================================================
\section{The Philosophy of Consequence}

\subsection{Narrative First}
\textbf{Every roll tells a story.} In Fate's Edge, mechanics serve the story, not the other way around. When you roll dice, you are not asking ``did I succeed?''—you are asking ``what happens next?''

\begin{sidebar}
    \textbf{The Golden Rule:} When in doubt, make the choice that serves the story. Set DV=3, Position=Controlled, and let the dice fall.
\end{sidebar}

\subsection{Safety at the Table}
Fate's Edge explores mature themes. Use these tools to keep play safe:
\begin{itemize}
    \item \textbf{X-Card:} Any player may say ``X'' to stop a scene. No questions asked.
    \item \textbf{Lines:} Content that will not appear (e.g., ``no harm to children'').
    \item \textbf{Veils:} Content that fades to black (e.g., ``romance is implied, not described'').
\end{itemize}
Establish these in Session Zero.

%=====================================================================
% SECTION 2: CORE RESOLUTION
%=====================================================================
\section{The Core Resolution}

\subsection{The 90-Second Loop}
Every meaningful action follows this cycle:
\begin{enumerate}
    \item \textbf{Declare} your intent and approach (Attribute + Skill).
    \item \textbf{GM sets DV} (2–5+) and \textbf{Position} (Dominant/Controlled/Desperate).
    \item \textbf{Roll} dice pool (Attribute + Skill d10s).
    \item \textbf{Count successes} (6–10 = 1 success). Each 1 = Story Beat (SB) for GM.
    \item \textbf{Apply outcome} using the Outcome Matrix.
    \item \textbf{GM spends SB} to introduce complications.
\end{enumerate}

\subsection{Difficulty Value (DV)}
\begin{coretable}[Difficulty Values]
\begin{tabular}{>{\bfseries}p{0.25\linewidth} p{0.65\linewidth}}
\toprule
\textbf{DV} & \textbf{When to Use} \\
\midrule
2 & Routine – clear path, no pressure \\
3 & Pressured – mild opposition (default) \\
4 & Hard – serious resistance \\
5+ & Extreme – dramatic gamble \\
\bottomrule
\end{tabular}
\end{coretable}
\textbf{Default DV is 3.} Do not inflate DVs to ``challenge'' players—use Position instead.

\subsection{Position}
Position tells you how risky the action is and changes how dice behave:
\begin{coretable}[Position Effects]
\begin{tabular}{>{\bfseries}p{0.3\linewidth} >{\bfseries}p{0.35\linewidth} >{\bfseries}p{0.25\linewidth}}
\toprule
\textbf{Position} & \textbf{Frame} & \textbf{Mechanical Effect} \\
\midrule
Dominant & You press advantage & Re-roll one failure \\
Controlled & Balanced norm & No re-rolls \\
Desperate & Act under duress & Re-roll one success \\
\bottomrule
\end{tabular}
\end{coretable}
\textbf{Pro tip:} Spend 1 Boon to improve Position by one step before rolling.

\subsection{The Outcome Matrix}
\begin{coretable}[Outcome Matrix]
\begin{tabular}{>{\bfseries}p{0.3\linewidth} p{0.6\linewidth}}
\toprule
\textbf{Roll Result} & \textbf{What Happens} \\
\midrule
Successes $\ge$ DV, SB = 0 & \textbf{Clean Success –} Intent achieved, no complication. \\
Successes $\ge$ DV, SB > 0 & \textbf{Success with SB –} Intent achieved; GM spends SB. \\
$0 <$ Successes $<$ DV & \textbf{Partial –} Progress; gain 1 Boon. \\
Successes = 0 & \textbf{Miss –} No progress; gain 2 Boons; GM spends SB. \\
\bottomrule
\end{tabular}
\end{coretable}

\textbf{A Partial is not failure.} It moves the story forward: improve Position for next attempt, reduce DV, succeed at reduced scale, or reveal information.

\textbf{A Miss is never ``nothing happens.''} The GM must introduce a complication.

\subsection{Critical Success (10s)}
A roll of 10 counts as \textbf{two successes}. If the overall roll succeeds:
\begin{itemize}
    \item \textbf{1 ten:} Strong success—improve Position or gain minor advantage.
    \item \textbf{2 tens:} Exceptional—choose two benefits.
    \item \textbf{3+ tens:} Legendary—decisive victory; clear 1 segment from a relevant clock.
\end{itemize}

%=====================================================================
% SECTION 3: BOONS AND STORY BEATS
%=====================================================================
\section{Resources: Boons \& Story Beats}

\subsection{Boons—Player Currency}
Boons represent luck, momentum, and narrative leverage. \textbf{Spend them freely.}
\begin{coretable}[Boons at a Glance]
\begin{tabular}{>{\bfseries}p{0.3\linewidth} >{\bfseries}p{0.15\linewidth} p{0.45\linewidth}}
\toprule
\textbf{Earning Boons} & \textbf{Gain} & \textbf{Spending Boons} \\
\midrule
Partial success & Gain 1 & Re-roll a die (1 Boon) \\
Miss & Gain 2 & Improve Position (1 Boon) \\
Bond-driven action & Gain 1 (once/bond/session) & Activate Asset (1 Boon) \\
GM award & +1 (rare) & Convert 2 $\to$ 1 XP (once/session) \\
\bottomrule
\end{tabular}
\end{coretable}
\textbf{Holding cap:} 5 Boons. At end of scene, reduce to 2 (excess lost).

\begin{sidebar}
    \textbf{The Math of Fun:} Hoarding Boons for XP conversion starves the table of drama. A re-rolled failure or an improved Position creates a memorable moment. Spend Boons!
\end{sidebar}

\subsection{Story Beats (SB)—GM Currency}
Every time a player rolls a 1, the GM gains Story Beats. Spend SB to introduce complications:
\begin{coretable}[SB Spend Menu]
\begin{tabular}{>{\bfseries}p{0.2\linewidth} p{0.7\linewidth}}
\toprule
\textbf{Cost} & \textbf{Complication} \\
\midrule
1 SB & Minor – noise, distraction, tick a clock +1 \\
2 SB & Moderate – alarm raised, lose advantage, worsen Position \\
3+ SB & Major – reinforcements, scene shifts, break an asset \\
\bottomrule
\end{tabular}
\end{coretable}
\textbf{Be aggressive with SB.} Unused pressure is wasted momentum.

%=====================================================================
% SECTION 4: HARM AND FATIGUE
%=====================================================================
\section{Injury \& Exhaustion}

\subsection{Fatigue—The Weight of the World}
Each point of Fatigue worsens your Position. If already Desperate, each Fatigue imposes $-1$ die instead.
\begin{coretable}[Fatigue Progression]
\begin{tabular}{>{\bfseries}p{0.35\linewidth} p{0.55\linewidth}}
\toprule
\textbf{Fatigue Level} & \textbf{Effect} \\
\midrule
1 & Position worsens one step \\
2 & If already Desperate, $-1$ die; otherwise Position worsens \\
3 & $-2$ dice (if already Desperate from Fatigue) \\
4+ & Collapse, KO, incapacitated \\
\bottomrule
\end{tabular}
\end{coretable}
\textbf{Recovery:} Short rest removes 2 Fatigue; full night's rest removes all.

\subsection{Harm—Serious Injury}
\begin{coretable}[Harm Levels]
\begin{tabular}{>{\bfseries}p{0.3\linewidth} p{0.6\linewidth}}
\toprule
\textbf{Harm Level} & \textbf{Effect} \\
\midrule
Harm 1 & $-1$ die on related actions \\
Harm 2 & $-1$ dice on most actions \\
Harm 3 & Incapacitated or dying \\
\bottomrule
\end{tabular}
\end{coretable}

\subsection{The Roller-Coaster Effect}
\textbf{Taking Harm clears Fatigue.} The shock of a wound resets your exhaustion. This creates dramatic swings—a character near collapse who takes a wound and finds a second wind.

\subsection{Armor Conversion}
When you take Harm, armor converts some to Fatigue:
\begin{coretable}[Armor Conversion]
\begin{tabular}{>{\bfseries}p{0.25\linewidth} >{\bfseries}p{0.25\linewidth} >{\bfseries}p{0.25\linewidth} >{\bfseries}p{0.25\linewidth}}
\toprule
\textbf{Armor} & \textbf{Harm 1} & \textbf{Harm 2} & \textbf{Harm 3} \\
\midrule
Light & Fatigue 1 & Harm 1 + Fatigue 1 & Harm 2 + Fatigue 1 \\
Medium & Fatigue 1 & Fatigue 1 & Harm 2 + Fatigue 1 \\
Heavy & Fatigue 1 & Fatigue 1 & Fatigue 2 \\
\bottomrule
\end{tabular}
\end{coretable}

%=====================================================================
% SECTION 5: CHARACTER CREATION
%=====================================================================
\section{Building Your Character}

\subsection{10-Minute Checklist}
Start with \textbf{30 XP} (max 34 with Bonds/Complications).
\begin{enumerate}[leftmargin=*]
    \item \textbf{Concept:} One sentence (origin, profession, defining trait).
    \item \textbf{Attributes (1–5):} Spend XP—primary to 3 (9 XP), secondaries to 2 (6 XP each).
    \item \textbf{Skills (0–5):} Spend 8–12 XP on key skills.
    \item \textbf{Talents:} 1–3 talents (2–6 XP each).
    \item \textbf{Patron (optional):} Choose a Patron for magic.
    \item \textbf{Assets/Followers (optional):} Minor Asset (4 XP) or follower (Cap² XP).
    \item \textbf{Bonds \& Complications:} Up to +4 XP total.
\end{enumerate}

\subsection{Attributes}
\begin{coretable}[Attributes]
\begin{tabular}{>{\bfseries}p{0.25\linewidth} >{\bfseries}p{0.45\linewidth} >{\bfseries}p{0.2\linewidth}}
\toprule
\textbf{Attribute} & \textbf{What It Does} & \textbf{Example Skills} \\
\midrule
Body & Strength, agility, endurance & Melee, Athletics, Stealth \\
Wits & Perception, quick thinking & Insight, Lore, Investigation \\
Spirit & Willpower, intuition & Arcana, Resolve, Survival \\
Presence & Charisma, leadership & Sway, Command, Deception \\
\bottomrule
\end{tabular}
\end{coretable}
\textbf{Cost to raise:} New rating $\times$ 3 XP. Downtime = new rating in days.

\subsection{Skills}
\textbf{Cost to raise:} New level $\times$ 2 XP. Downtime = new level in days.
\begin{coretable}[Core Skills]
\begin{tabular}{>{\bfseries}p{0.25\linewidth} >{\bfseries}p{0.25\linewidth} >{\bfseries}p{0.4\linewidth}}
\toprule
\textbf{Combat} & \textbf{Physical} & \textbf{Social} \\
\midrule
Melee & Athletics & Sway \\
Ranged & Stealth & Command \\
Brawl & Endurance & Deception \\
Tactics & Craft & Performance \\
& Survival & Insight \\
\addlinespace
\textbf{Knowledge} & \textbf{Magic} & \\
Lore & Arcana & \\
Investigation & Ritual & \\
Medicine & & \\
\bottomrule
\end{tabular}
\end{coretable}

\subsection{Sample Starting Build—The Duelist (30 XP)}
\begin{itemize}[leftmargin=*]
    \item \textbf{Attributes:} Body 3 (15 XP), Wits 2 (6 XP)
    \item \textbf{Skills:} Melee 2 (6 XP), Athletics 1 (2 XP)
    \item \textbf{Talent:} Keen Senses (2 XP)
    \item \textbf{Total:} 31 XP (bank 1)
\end{itemize}

%=====================================================================
% SECTION 6: THE REGION—ACASIA
%=====================================================================
\section{The Known World: Acasia}
\textit{``The Broken Province''}

\subsection{Genre \& Mood}
\textbf{Low fantasy / mercenary realism / border feud.} Acasia is a former province of the Utaran Empire that festers in the space between famine and ambition. Every hill wears a crown, every bridge demands a price, and every road curves back to the same cursed crossroads.

\subsection{What You Notice First}
The smell of mildew on old parchment and wet ash. The way every milestone has been recarved. The silence of a village that heard you coming and chose not to open its gates.

\subsection{The One Rule That Kills Newcomers}
\textbf{Never sign a contract on a bridge.} The river is a witness, and it never forgets a signature.

\subsection{Key Factions}
\begin{coretable}[Acasia Factions]
\setlength{\tabcolsep}{1.5pt} % Local override for optimal spacing
\begin{tabular}{>{\bfseries}p{0.28\linewidth} >{\bfseries}p{0.22\linewidth} >{\bfseries}p{0.4\linewidth}}
\toprule
\textbf{Faction} & \textbf{Leader} & \textbf{Goal} \\
\midrule
Broken March & Margravine Mira & Keep province from fracturing \\
Lame King's Claimants & Oren ``Tin Crown'' & Be recognised as true sovereign \\
Free Companies & Various captains & Find permanent contracts \\
Blackwood Matriarchs & Hedge-witches & Survive the Curse \\
\bottomrule
\end{tabular}
\end{coretable}

\subsection{The Curse of Acasia}
The land remembers what we did—and what we failed to do. Every broken milestone is a forgotten promise. Every abandoned vineyard is a harvest we chose not to plant.
\textbf{Mechanical effect:} When crossing a bridge or crossroads, roll Spirit + Resolve (DV 3). On a failure, the Curse stirs—gain 1 Fatigue or a minor complication (GM's choice).

\subsection{Acasia Mini-Generator (d6)}
\begin{enumerate}[leftmargin=*]
    \item A bridge toll has tripled overnight; the toll-keeper is a deserter from your old company.
    \item A hedge-witch offers to break your blood‑debt—for the memory of your first kill.
    \item The Lame King travels with a mute jester who writes prophecies on fallen leaves.
    \item Bandits stole a Blackwood Matriarch's curse‑box. Recover it before it opens.
    \item A free company captain wants to hire you—but their last three hires died of ``accidents.''
    \item The Margravine's seal is missing. Whoever finds it can rewrite border law.
\end{enumerate}

%=====================================================================
% SECTION 7: RUNNING THE GAME (GM)
%=====================================================================
\section{Running Fate's Edge}

\subsection{The GM Mindset}
You are not a referee or a narrator robot. You are a \textbf{co-creator of drama}. Your tools (SB, clocks) are not punishments—they are how you play the game. When you spend a Story Beat to tick a clock, you are making a move. Enjoy it.

\subsection{Your Defaults}
When in doubt:
\begin{itemize}[leftmargin=*]
    \item \textbf{DV: 3}
    \item \textbf{Position: Controlled}
    \item \textbf{Consequence:} Fatigue, time pressure, or exposure
    \item \textbf{Clocks:} One short clock tied to the scene
\end{itemize}

\subsection{What You Are Doing}
\begin{itemize}[leftmargin=*]
    \item Stating Position clearly and confidently
    \item Letting partials move the story forward
    \item Spending Story Beats to make failure interesting
    \item Allowing success to create new problems
\end{itemize}

\subsection{The Three-Clock Guideline}
Keep 1–3 active clocks per scene:
\begin{coretable}[Clock Sizes]
\begin{tabular}{>{\bfseries}p{0.25\linewidth} p{0.65\linewidth}}
\toprule
\textbf{Segments} & \textbf{When to Use} \\
\midrule
4 & Short, straightforward challenge \\
6 & Standard scene or tension \\
8 & Extended or complex problem \\
10 & Epic or campaign-long pressure \\
\bottomrule
\end{tabular}
\end{coretable}

\subsection{GM Survival Tips}
\subsubsection{When Players Go Off the Rails}
Do not panic. Ask: \textit{``What is your goal?''} Name the simplest obstacle—then set Position and DV.

\subsubsection{When Players Hoard Boons}
Remind them: \textit{``Boons are not XP tokens—they are dramatic fuel.''} Introduce time pressure that forces spending.

\subsubsection{When a Scene Stalls}
\begin{itemize}[leftmargin=*]
    \item Ask a leading question: \textit{``What are you afraid of here?''}
    \item Tick a clock: \textit{``While you argue, the fire spreads—Collapse [4] advances 1.''}
    \item Cut to another player.
\end{itemize}

\subsubsection{When You Don't Know a Rule}
Make a ruling that keeps the story moving. Default to DV 3, Controlled Position. Look it up later.

%=====================================================================
% SECTION 8: THE STARTER ADVENTURE
%=====================================================================
\section{The Lantern at Dusk}
\textit{A Starter Adventure in Acasia}

\subsection{Setup}
The party is hired by \textbf{Elder Sarai} of the remote village of Duskwood. A terrible blight is spreading from the old barrow in the hills—the same barrow where a cursed lantern was buried centuries ago. The lantern must be brought to a stone circle in the valley to break the curse.

\subsection{What the Party Knows}
\begin{itemize}[leftmargin=*]
    \item The lantern is in the central chamber of the barrow.
    \item The barrow is guarded by restless spirits (the Duskwights).
    \item A rival treasure hunter, \textbf{Quick Lena}, has also entered the barrow.
    \item The curse will kill the village's crops in three days.
\end{itemize}

\subsection{Clocks}
\begin{coretable}[Adventure Clocks]
\begin{tabular}{>{\bfseries}p{0.3\linewidth} p{0.6\linewidth}}
\toprule
\textbf{Clock} & \textbf{Effect} \\
\midrule
Barrow Collapse [6] & Ticks on failure inside. At 6, entrance seals. \\
Lena's Agenda [4] & Ticks each scene inside. At 4, Lena takes lantern first. \\
\bottomrule
\end{tabular}
\end{coretable}

\subsection{Scenes}
\subsubsection{Scene 1—The Entry (DV 3, Controlled)}
Obstacle: Cave-in blocks the way (Body + Athletics). \textbf{Partial:} Tick Barrow Collapse +1. \textbf{Miss:} Mark Harm 1, tick Collapse +1.

\subsubsection{Scene 2—The Spirit Corridor (DV 3, Controlled)}
Three Duskwights demand an offering (a memory or a secret). Roll Presence + Sway. \textbf{Partial:} They also steal random equipment. \textbf{Miss:} They attack (Cap 2, Harm 1 each).

\subsubsection{Scene 3—The Lantern Chamber}
Quick Lena is trying to pry the lantern loose. Confront her: Presence + Sway (DV 3, Controlled). \textbf{Partial:} She agrees to split the reward but demands a favor later (start Debt Clock [4]). \textbf{Miss:} She triggers a trap—tick Collapse +2.

\subsubsection{Scene 4—Escape (Skill Challenge)}
The party needs \textbf{3 successes before 2 failures} to escape. Each PC describes their approach. DV 3, Position Desperate. Each failure ticks Collapse +1.

\subsection{Resolution}
Return the lantern to the stone circle. Break the curse with Spirit + Lore (DV 2). \textbf{Partial:} Curse lifts, but lantern attracts another danger soon. \textbf{Miss:} Curse pushed back one year; a dark spirit escapes.

\subsection{Advancement}
Each PC earns \textbf{6–10 XP} at the end.

%=====================================================================
% SECTION 9: MAGIC (BASIC)
%=====================================================================
\section{Basic Magic}

\subsection{The Five Paths (Quick Reference)}
\begin{coretable}[Magic Paths]
\begin{tabular}{>{\bfseries}p{0.3\linewidth} >{\bfseries}p{0.2\linewidth} >{\bfseries}p{0.25\linewidth} >{\bfseries}p{0.2\linewidth}}
\toprule
\textbf{Path} & \textbf{Complexity} & \textbf{Cost Type} & \textbf{Best For} \\
\midrule
Runkeeper & Low & Obligation & First-timers, reliable casters \\
Cantor & Medium & Corruption & Bards, risk-takers \\
Invoker & Medium & Obligation + time & Ritual specialists \\
Summoner & High & Leash + Boons & Pet-class fans \\
Free Caster & High & Backlash & Improv players \\
\bottomrule
\end{tabular}
\end{coretable}
\textbf{For first magical characters, play a Runkeeper.}

\subsection{Runkeeper—The Reliable Path}
\textbf{Requirements:} Familiar (Thiasos, 2 XP) + Codex (4 XP) + Patron.

\textbf{Invoking a Rite:} 1 action. Mark +1 Obligation. Effect works as written.

\textbf{Push It:} Once per scene, mark +1 Obligation to amplify (+1 die or +1 Effect).

\textbf{Obligation Capacity:} Spirit + Presence. Exceeding capacity: 1 Fatigue per segment.

\subsection{Sample—The Traveler's Runkeeper}
\begin{sidebar}
    \textbf{Kestra, Traveler's Runkeeper} (34 XP)
    \begin{itemize}[leftmargin=*]
        \item Spirit 3 (9 XP), Wits 2 (6 XP)
        \item Lore 2 (6 XP), Insight 1 (2 XP), Sway 1 (2 XP)
        \item Talents: Familiar (2 XP), Codex (4 XP)
        \item Rites: Road-Sense (4 XP), Traveler's Boon (5 XP)
        \item Obligation Capacity: 4 (Spirit 3 + Presence 1)
    \end{itemize}
\end{sidebar}

\subsection{Free Casting—The TAGS System}
For improvisational magic, construct spells using TAGS:
\[
\text{DV} = 1 + \text{number of TAGS}
\]
Dangerous TAGS (Leap, Fold, Gate, Transform, Dominate) add $+2$ DV.

\subsubsection{Example Spells}
\begin{itemize}[leftmargin=*]
    \item \textbf{Light [Illuminate]} – DV 2. Create small, sustained light.
    \item \textbf{HEAL [HEAL]} – DV 2. Remove 1 Fatigue from self or touched ally.
    \item \textbf{Shield of Air [WARD] [DEFEND]} – DV 3. Gain Position +1 on next defense.
    \item \textbf{Flame Cloak [BURNING] [AREA]} – DV 3. Enemies suffer $-1$ die to approach.
\end{itemize}
\textbf{Backlash:} On a Miss or when GM spends SB, magic backfires—Fatigue 1, Harm 1, or a thematic complication.

%=====================================================================
% SECTION 10: ADVANCEMENT
%=====================================================================
\section{Advancement \& XP}

\subsection{Earning XP (Per Session)}
\begin{coretable}[XP Breakdown]
\begin{tabular}{>{\bfseries}p{0.5\linewidth} >{\bfseries}p{0.4\linewidth}}
\toprule
\textbf{Achievement} & \textbf{XP} \\
\midrule
Attendance & +2 \\
Major objectives & +2–4 \\
Discoveries & +1–2 \\
Hard choices & +1–2 \\
Story engagement & +1–3 \\
Personal goals & +1–2 \\
\bottomrule
\end{tabular}
\end{coretable}
\textbf{Game pace options:} Gritty (4–6 XP/session), Standard (6–10), Epic (10–14).

\subsection{Spending XP}
\begin{coretable}[XP Costs]
\begin{tabular}{>{\bfseries}p{0.4\linewidth} >{\bfseries}p{0.25\linewidth} >{\bfseries}p{0.25\linewidth}}
\toprule
\textbf{Improvement} & \textbf{XP Cost} & \textbf{Downtime} \\
\midrule
Attribute increase & New rating $\times$ 3 & New rating in days \\
Skill increase & New level $\times$ 2 & New level in days \\
Minor Asset & 4 XP & 1 day \\
Standard Asset & 8 XP & 1 week \\
Major Asset & 12 XP & 1 month \\
Minor Talent & 2–4 XP & 3–10 days \\
Major Talent & 5–8 XP & 1–2 weeks \\
\bottomrule
\end{tabular}
\end{coretable}

\subsection{Tiers of Play}
\begin{coretable}[Character Tiers]
\begin{tabular}{>{\bfseries}p{0.25\linewidth} >{\bfseries}p{0.4\linewidth} >{\bfseries}p{0.25\linewidth}}
\toprule
\textbf{Tier} & \textbf{XP Range} & \textbf{Reputation} \\
\midrule
I: Novice & 0–40 & Local \\
II: Seasoned & 41–90 & Regional \\
III: Veteran & 91–150 & National \\
IV: Paragon & 151–220 & Continental \\
V: Mythic & 221+ & Legendary \\
\bottomrule
\end{tabular}
\end{coretable}

%=====================================================================
% SECTION 11: QUICK REFERENCE
%=====================================================================
\section{Quick Reference Tables}

\subsection{Combat in Brief}
\begin{itemize}[leftmargin=*]
    \item \textbf{Turn:} 1 Action + 1 Move
    \item \textbf{Move:} Shift one range band (Close $\leftrightarrow$ Near $\leftrightarrow$ Far)
    \item \textbf{Opportunity Attack:} Leaving Close range without Disengage provokes one free melee attack.
    \item \textbf{Sidestep Strike:} Move one band as part of melee attack; worsen Position by one step.
\end{itemize}

\subsection{Weapon Modifiers}
\begin{coretable}[Melee Weapons]
\begin{tabular}{>{\bfseries}p{0.25\linewidth} >{\bfseries}p{0.25\linewidth} >{\bfseries}p{0.4\linewidth}}
\toprule
\textbf{Weight} & \textbf{Close} & \textbf{Near} & \textbf{Notes} \\
\midrule
Light & +2d & +1d & Quick in tight quarters \\
Medium & +1d & +2d & Set once/scene or $-1$d first attack \\
Heavy & $-1$d & +3d & Set once/scene or $-2$d first attack \\
\bottomrule
\end{tabular}
\end{coretable}

\subsection{Range Bands}
\begin{itemize}[leftmargin=*]
    \item \textbf{Close:} Touching distance—grapples, knife-work
    \item \textbf{Near:} Same room/area—a rush away
    \item \textbf{Far:} Same site but distant—requires route or time
    \item \textbf{Absent:} Off-screen—requires scene change
\end{itemize}

\subsection{Skill Synergies—Creative Pairs}
\begin{coretable}[Alternative Skill Pairs]
\begin{tabular}{>{\bfseries}p{0.25\linewidth} >{\bfseries}p{0.35\linewidth} >{\bfseries}p{0.3\linewidth}}
\toprule
\textbf{Action} & \textbf{Standard Pair} & \textbf{Alternative Pair} \\
\midrule
Climbing & Body + Athletics & Wits + Athletics (clever route) \\
Persuading & Presence + Sway & Wits + Sway (logical argument) \\
Investigating & Wits + Investigation & Spirit + Investigation (intuitive leap) \\
Feinting & Wits + Melee & Presence + Melee (intimidation) \\
\bottomrule
\end{tabular}
\end{coretable}

\subsection{Final GM Mantra}
\begin{sidebar}
    \textbf{Set the DV low. Make the world react honestly. Let the pressure do the work. And remember—you are playing the game too.}
\end{sidebar}

%=====================================================================
% SECTION 12: SESSION ZERO CHECKLIST
%=====================================================================
\section{Session Zero Checklist}
Use this checklist to start your campaign.
\begin{enumerate}[leftmargin=*]
    \item \textbf{Lines \& Veils} – What content is off-limits or only implied?
    \item \textbf{Tone} – Cinematic action? Grim survival? Political intrigue? Dark horror?
    \item \textbf{Campaign length} – One-shot, short arc (3–6 sessions), or long campaign (12+)?
    \item \textbf{Character hooks} – Each player shares one goal and one fear.
    \item \textbf{Patron interest} – Which Patrons (if any) are already watching the party?
    \item \textbf{Starting situation} – Where does the story begin?
    \item \textbf{House rules} – Any optional subsystems you plan to use?
    \item \textbf{Scheduling \& attendance} – How often do you meet? What if a player misses?
    \item \textbf{Player-GM contract} – ``We are co-creators, not adversaries.''
    \item \textbf{Resource mindset} – ``Boons are not XP tokens. Spend them. Embrace failure.''
\end{enumerate}
\textit{``Now close the book—or keep reading. The road is waiting.''}

%=====================================================================
% BACK MATTER
%=====================================================================
\newpage
\thispagestyle{empty}
\vspace*{2cm}
\begin{center}
    {\LARGE\sffamily\itshape ``The coin that never spends is the one you don't remember taking.''}\\[1.5em]
    — Serafine of the Velvet Touch\\[2em]
    \rule{0.5\textwidth}{0.5pt}\\[1.5em]
    {\small Fate's Edge Essential Guide — CC BY 4.0 — Nicholas A. Gasper}\\[0.5em]
    {\tiny \textit{Set down in Castra Ferrum, Year of the Iron Harvest}}
\end{center}
\vspace*{1.5cm}

\end{document}
